% Short derivation of the model in lidar_eba.py, with the reasons for each parameter.
\documentclass[11pt,a4paper]{article}
\usepackage[margin=2.2cm]{geometry}
\usepackage{amsmath,amssymb}
\usepackage{booktabs,tabularx}
\usepackage[hidelinks]{hyperref}
\usepackage{parskip}

\title{InnovizTwo for Emergency Brake Assist:\\ derivation of the statistical model}
\author{3DSSF -- Kozma, Exercise 3}
\date{}

\begin{document}
\maketitle

\noindent\textbf{Question.} An object of width $W$, height $H$ and reflectivity $\rho$ is at distance $R$.
What is the probability that the lidar gets at least $k_{\min}=5$ returns from it in one frame,
and up to which $R$ is that probability at least 99\,\%?

\section{Number of beams on the object}

The lidar samples a regular angular grid with spacing $\Delta\theta_h$ (horizontal) and
$\Delta\theta_v$ (vertical). At distance $R$ neighbouring beams are $s = R\,\Delta\theta$ apart
(small angles: $\tan\Delta\theta \approx \Delta\theta$, error $< 10^{-6}$ for $0.1^\circ$).
One side of the object therefore spans
\[
  L = \frac{W}{R\,\Delta\theta_h} \quad\text{beam spacings.}
\]
\textbf{Random grid offset.} Beams sit at $x_j = u + j s$, with the offset $u$ uniform on $[0,s)$
(the object can be anywhere relative to the grid). The interval $[0,W]$ contains
$\lfloor L\rfloor + 1$ beams if $u < W - \lfloor L\rfloor s$, otherwise $\lfloor L\rfloor$. Hence
\[
  P\big(n_h = \lfloor L\rfloor + 1\big) = \frac{W - \lfloor L\rfloor s}{s} = L - \lfloor L\rfloor = \{L\},
  \qquad P\big(n_h = \lfloor L\rfloor\big) = 1 - \{L\}.
\]
The same holds vertically, independently. So $n = n_h n_v$ takes 4 values with known probabilities
$w_{ij}$. As a check, $\mathbb{E}[n_h] = L$: on average the count is exactly the geometric value.

\section{Detection probability of one beam}

\textbf{Signal.} For a Lambertian target larger than the laser spot, the lidar equation gives the
received power
\[
  P_r \propto \frac{\rho}{R^2}\, e^{-2\alpha R},
\]
where $e^{-2\alpha R}$ is the loss in the air on the way out and back (Beer--Lambert, extinction
coefficient $\alpha$).

\textbf{From signal to probability.} We model the detected signal events in one beam as Poisson
with a mean $\lambda$ proportional to the received power, $\lambda = C\rho R^{-2}e^{-2\alpha R}$.
A beam gives a return if at least one event is detected:
\[
  p(R,\rho) = 1 - e^{-\lambda} = 1 - \exp\!\Big(-C\,\frac{\rho}{R^2}\,e^{-2\alpha R}\Big).
\]
This is the simplest curve that is $0$ for no signal, tends to $1$ for a strong signal and has
only one free constant. That matters, because the datasheet gives only one point.

\textbf{Calibration.} The datasheet states $p(R_0,\rho_0) = P_0$ in clear air ($\alpha=0$):
\[
  1 - e^{-C\rho_0/R_0^2} = P_0 \;\Longrightarrow\; C = -\ln(1-P_0)\,\frac{R_0^2}{\rho_0}.
\]
With $R_0 = 210$\,m, $\rho_0 = 0.1$, $P_0 = 0.9$ this gives
$C\rho/R^2 = \ln 10 \cdot (R_0/R)^2$ for $\rho = 0.1$, so at $R=250$\,m:
$p = 1 - e^{-2.303\cdot 0.7056} = 1 - e^{-1.625} = \mathbf{0.803}$.

\section{Number of returns and the requirement}

Given $n$ beams on the object, each returning independently with probability $p$,
\[
  K \mid n \sim \mathrm{Binomial}(n,p), \qquad
  P(K \ge k \mid n) = 1 - \sum_{i=0}^{k-1} \binom{n}{i} p^i (1-p)^{n-i}.
\]
Averaging over the four beam counts (law of total probability):
\[
  P(K \ge k) = \sum_{i,j\in\{0,1\}} w_{ij}\, P(K \ge k \mid n_{ij}).
\]
\textbf{Verdict.} $R_{99} = \max\{R : P(K\ge 5) \ge 0.99\}$; the object passes if $R_{99} \ge 250$\,m.

\textbf{Worked example} (dark pedestrian, $0.5\times1.8$\,m, 250\,m, clear).
$s_h = 0.436$\,m, $L_h = 1.15$; $s_v = 0.218$\,m, $L_v = 8.25$.
\begin{center}
\begin{tabular}{cccc}
  \toprule
  $n$ & $w$ & $P(K\ge5\mid n)$ & contribution \\
  \midrule
  $1\times 8 = 8$  & $0.85\cdot0.75 = 0.64$ & 0.947 & 0.604 \\
  $1\times 9 = 9$  & $0.85\cdot0.25 = 0.21$ & 0.982 & 0.209 \\
  $2\times 8 = 16$ & $0.15\cdot0.75 = 0.11$ & 1.000 & 0.112 \\
  $2\times 9 = 18$ & $0.15\cdot0.25 = 0.04$ & 1.000 & 0.037 \\
  \midrule
  & & & $\mathbf{0.962} < 0.99$ \\
  \bottomrule
\end{tabular}
\end{center}
The mean is $\mathbb{E}[K] = \mathbb{E}[n]\,p = 1.15\cdot 8.25\cdot 0.803 = 7.6 \ge 5$:
the average passes, the probability fails.

\section{Weather}

Weather reports give the visibility $V$ (meteorological optical range), not $\alpha$. Visibility
is defined as the distance at which the contrast of a dark object drops to 2\,\%, so
(Koschmieder)
\[
  e^{-\alpha V} = 0.02 \;\Longrightarrow\; \alpha = \frac{-\ln 0.02}{V} = \frac{3.912}{V}.
\]
For each condition $c$ we evaluate the model with $\alpha_c$. The share of driving time $\pi_c$
gives the yearly figure, again by total probability:
\[
  P_{\text{weighted}}(K\ge5) = \sum_c \pi_c\, P(K\ge5 \mid \alpha_c)
  = 0.88\,P_{\text{clear}} + 0.08\,P_{\text{rain}} + 0.04\,P_{\text{fog}}.
\]
Pedestrian at 250\,m: $0.88\cdot0.962 + 0.08\cdot0.803 + 0.04\cdot0.047 = 0.913$.

\section{Parameters and why we chose them}

\small
\renewcommand{\arraystretch}{1.25}
\begin{tabularx}{\textwidth}{l l X}
  \toprule
  Parameter & Value & Reason / source \\
  \midrule
  \multicolumn{3}{l}{\textit{Lidar: InnovizTwo datasheet v3.1 (2024), Condor configuration, region of interest}} \\
  $\Delta\theta_h \times \Delta\theta_v$ & $0.1^\circ \times 0.05^\circ$ & datasheet, ROI resolution. The ROI ($30^\circ\times9.6^\circ$)
    is about $130\times40$\,m at 250\,m, so the lane ahead is inside it. \\
  $R_0,\ \rho_0,\ P_0$ & 210\,m, 10\,\%, 0.90 & datasheet: ``210\,m @ 10\,\% reflectivity'', note 4:
    ``True Positives = 90\,\% per pixel'' (100\,klx sunlight). This is exactly the per-beam probability the model needs. \\
  $R_{\max}$, frame rate & 300\,m, 20\,FPS & datasheet; results are capped at 300\,m. \\
  \midrule
  \multicolumn{3}{l}{\textit{Requirement and scenario}} \\
  $R_{\text{req}},\ k_{\min}$ & 250\,m, 5 & given by the task. \\
  confidence & 99\,\% & our choice: an EBA must not miss an object in more than about 1 frame in 100.
    Sensitivity: 95\,\% gives 253\,m for the pedestrian, 99.9\,\% gives 219\,m. \\
  $\rho$ & 10\,\% & the standard worst case in lidar datasheets; dark clothing and tyres are 5--10\,\%.
    50\,\% (bright clothing) is shown for comparison. \\
  pedestrian & $0.5\times1.8$\,m & adult shoulder width $\approx 0.5$\,m, height $\approx 1.75$--$1.8$\,m. \\
  car & $1.8\times1.5$\,m & rear of a typical passenger car (width 1.75--1.85\,m, height 1.45--1.5\,m). \\
  motorcycle & $0.8\times1.4$\,m & motorcycle with rider seen from behind (handlebar/mirrors, rider's head). \\
  truck & $2.5\times3.5$\,m & rear of a truck/trailer; EU limits are 2.55\,m wide and 4\,m high. \\
  tyre / cargo & $0.5\times0.5$\,m & lost load or a car tyre lying on the road (tyre diameter 0.6--0.7\,m, seen partly). \\
  \midrule
  \multicolumn{3}{l}{\textit{Weather}} \\
  clear & $\alpha = 0$ & datasheet conditions; visibility $> 20$\,km gives $\alpha < 0.0002$/m, negligible at 250\,m. \\
  rain / haze & $V = 5$\,km & moderate rain or haze; $\alpha = 7.8\cdot10^{-4}$/m. \\
  fog & $V = 1$\,km & the WMO definition of fog is visibility below 1\,km, so this is the mildest fog; $\alpha = 3.9\cdot10^{-3}$/m.
    Denser fog is worse. \\
  shares $\pi_c$ & 88 / 8 / 4\,\% & our estimate, not official: Budapest has about 80 days per year with $\ge 1$\,mm rain,
    but rain lasts hours, not days; fog is mostly in winter. To be replaced by HungaroMet data. \\
  \midrule
  \multicolumn{3}{l}{\textit{Context: stopping distance}} \\
  speed & 130\,km/h & Hungarian motorway speed limit. \\
  latency, braking & 0.5\,s, 6\,m/s$^2$ & system reaction; firm but not full braking (dry road allows 8--9\,m/s$^2$).
    $d = vt + v^2/2a \approx 127$\,m, so 250\,m leaves 68 frames. \\
  \bottomrule
\end{tabularx}
\normalsize

Datasheet PDF: \url{https://levelfivesupplies.com/wp-content/uploads/2023/08/InnovizTwo-GMSL-Series-Production-Datasheet_v3.1-1.pdf}

\section{Main assumptions}
Flat rectangular targets facing the lidar; beams are points (no partial hits); beams fail
independently; the shape of $p(R)$ beyond the datasheet point follows from the Poisson model
(another common shape, the speckle model, gives 248\,m instead of 235\,m for the pedestrian, so ours
is the cautious one); false positives (1\,\% per pixel) are ignored; one frame is judged on its own.

\end{document}
